% !TEX program = xelatex
% ============================================================
% Overleaf 编译设置：
% 1. 点击左上角 Menu
% 2. 在 Compiler 下拉菜单中选择 "XeLaTeX"
% 3. 点击 Recompile
% ============================================================
\documentclass[sigconf,natbib=true]{acmart}
%% SIGIR 使用 ACM sigconf 格式（双栏布局）

%% 移除 ACM 特定元素以获得更简洁的评审版本
\settopmatter{printacmref=false}
\renewcommand\footnotetextcopyrightpermission[1]{}
\setcopyright{none}
\acmConference[SIGIR'25]{SIGIR 2025}{July 2025}{Padua, Italy}
\acmYear{2025}
\acmDOI{}
\acmISBN{}

\usepackage{booktabs}
\usepackage{multirow}
\usepackage{amsmath}
\let\Bbbk\relax  % 解决 amssymb 与其他包的冲突
\usepackage{amssymb}
\usepackage{algorithm}
\usepackage{algpseudocode}
\usepackage{graphicx}
\usepackage{subcaption}
\usepackage{xspace}
\usepackage{enumitem}
\usepackage{ifthen}
\usepackage{balance}

% XeLaTeX 中文支持 - 使用 Overleaf 预装的 Noto 字体
\usepackage{fontspec}
\usepackage{xeCJK}
\setCJKmainfont{Noto Serif CJK SC}
\setCJKsansfont{Noto Sans CJK SC}
\setCJKmonofont{Noto Sans Mono CJK SC}

\newcommand{\model}{\textsc{MV-Align}\xspace}
\newcommand{\base}{\textsc{SASRec}\xspace}
\newcommand{\tfidf}{\textsc{TF-IDF}\xspace}
\newcommand{\llm}{\textsc{LLM}\xspace}
\newcommand{\InfoNCE}{InfoNCE\xspace}

\begin{document}

\title{面向序列推荐的多视图文本特征：\\
显式交叉与对齐实现可靠的全排序增益}

% 评审版本：匿名作者
% 定稿示例：
% \author{第一作者}
% \affiliation{\institution{机构}\country{国家}}
% \email{email@example.com}

\begin{abstract}
文本特征已成为序列推荐器的标准插件，但两个实际问题仍不清楚：
(\emph{i}) 当轻量级文本特征（如 \tfidf）已经足够强时，多视图 \llm 特征能带来什么额外价值；以及
(\emph{ii}) 为什么在采样评估下的改进有时无法转化为全排序增益。
本文提出 \model，一个可工业部署的框架，结合了：(1) \emph{归一化多视图文本表示}（一个 \tfidf 基础视图加多个 \llm 提示视图），
(2) 一个\emph{显式交叉模块}，用于构建超越自注意力的高阶特征交互，以及
(3) \emph{多视图对比对齐}，将每个文本视图与协同ID空间绑定，提高全排序下的鲁棒性。
我们在两个亚马逊领域（美妆和玩具游戏）进行评估，并报告全排序下的标准和流行度分层指标。
在所有数据集上，\model 持续改善整体和长尾性能，同时在节约型（仅 \tfidf）和更强配置（多视图 \llm）之间提供清晰的成本收益权衡。
我们进一步分析采样指标与全排序指标之间的差距，并展示对齐如何缓解排序不稳定性。
\end{abstract}

\keywords{序列推荐, 文本特征, 多视图表示, 显式交叉, 对比对齐, 评估指标}

\maketitle

\section{引言}
诸如 \base~\cite{kang2018sasrec} 等序列推荐模型已成为建模用户行为序列的强大骨干网络。
同时，物品文本（标题、描述、评论）提供了补充的语义线索，能够实现更好的泛化，特别是对于冷启动和长尾物品。
近年来，\llm 衍生的表示进一步提升了文本质量，但也引发了关于特征成本、存储和推理延迟的实际关注。

本工作定位于务实的\emph{方法+实证}立场：
我们不是提出新的 \llm 范式，而是研究一个可部署的插件框架，该框架结合了多视图文本特征（来自 \tfidf 和 \llm 提示）、显式交叉交互和对比对齐。
我们的目标是回答：\emph{何时节约方案就足够了，何时多视图 \llm 文本值得投入？}

\paragraph{动机：采样指标 vs 全排序指标}
工业实践通常依赖采样评估（如均匀100个负样本）以提高速度，但采样指标可能与全排序不一致，可能导致模型排序错误~\cite{krichene2020sampledmetrics}。
因此，我们在主要结果中强调全排序评估，并提供采样与全排序差异的分析。

\section{研究问题}
我们围绕以下研究问题（RQs）组织研究：
\begin{itemize}[leftmargin=*,itemsep=2pt]
  \item \textbf{RQ1（多视图文本的价值）} 在全排序下，多视图 \llm 文本视图相比强 \tfidf 文本基线能改进多少？
  \item \textbf{RQ2（显式交叉的作用）} 显式交叉模块是否在 Transformer 骨干之外提供补充增益？
  \item \textbf{RQ3（对齐与评估稳定性）} 多视图对齐能否减少采样与全排序指标的差距并提高排序稳定性？
  \item \textbf{RQ4（成本收益）} 多视图 \llm 需要多少额外的计算/内存/延迟，在哪些场景下值得？
\end{itemize}

\section{方法}
\subsection{骨干网络：带文本融合的 SASRec}
设用户交互序列为 $s_u = [i_1,\dots,i_T]$，最大长度为 $L$。
\base 通过 Transformer 将 $s_u$ 编码为序列表示 $\mathbf{h}_u \in \mathbb{R}^d$。
每个物品 $i$ 有一个 ID 嵌入 $\mathbf{e}_i \in \mathbb{R}^d$。
我们用文本特征增强物品，并将它们融合成用于打分的\emph{融合物品嵌入} $\tilde{\mathbf{e}}_i$。

\subsection{多视图文本表示}
我们假设有一个轻量级 \tfidf 嵌入 $\mathbf{t}_i^{(0)}$ 作为基础视图，以及 $V$ 个 \llm 提示视图 $\{\mathbf{t}_i^{(v)}\}_{v=1}^V$（如身份/功能/受众/类别视图）。
每个视图被投影到骨干维度，并通过压缩激励（SENet）门控进行精炼：
\begin{align}
  \mathbf{z}_i^{(v)} &= \mathbf{W}_v \mathbf{t}_i^{(v)} \in \mathbb{R}^d,\\
  \mathbf{r}_i^{(v)} &= \mathbf{z}_i^{(v)} \odot \sigma(\mathrm{MLP}_v(\mathbf{z}_i^{(v)})).
\end{align}
在 \model 中我们应用\emph{逐视图} $\ell_2$ 归一化：
\begin{equation}
  \hat{\mathbf{r}}_i^{(v)} = \frac{\mathbf{r}_i^{(v)}}{\lVert \mathbf{r}_i^{(v)} \rVert_2}.
\end{equation}
我们为每个视图学习一个标量门控 $g_v=\sigma(\theta_v)\in(0,1)$，且\emph{不}强制跨视图归一化（每个视图独立贡献 $0$--$1$）。
所有视图（可选包括 \tfidf 基础视图）被拼接并投影回 $\mathbb{R}^d$：
\begin{equation}
  \mathbf{t}_i = \mathbf{P}\Big[\mathbf{t}_i^{(0)};\; g_1\hat{\mathbf{r}}_i^{(1)};\; \dots;\; g_V\hat{\mathbf{r}}_i^{(V)}\Big].
\end{equation}

\subsection{显式交叉融合}
我们不仅仅依赖自注意力来捕获协同信号和文本信号之间的高阶交互，而是引入基于 DCN-V2~\cite{wang2021dcnv2} 的\emph{显式交叉}模块。
给定拼接特征 $\mathbf{x}_i=[\mathbf{e}_i;\alpha_i \mathbf{t}_i]$，交叉层更新为
\begin{equation}
  \mathbf{x}^{(l+1)} = \mathbf{x}^{(l)} + \mathbf{x}^{(0)} \odot (\mathbf{W}_l\mathbf{x}^{(l)} + \mathbf{b}_l),
\end{equation}
并与浅层 MLP 结合产生融合嵌入 $\tilde{\mathbf{e}}_i$。
标量 $\alpha_i$ 包含全局文本门控和可选的逐物品门控；它还通过温度相关因子进行缩放以实现稳定融合。

\subsection{多视图对比对齐}
我们使用 \InfoNCE~\cite{oord2018cpc} 将每个文本视图与 ID 嵌入空间对齐。
对于一批正样本物品 $\{i\}_{i=1}^B$，定义相似度矩阵
$\mathbf{S}_{ij} = \mathrm{cos}(\mathbf{e}_{i}, \hat{\mathbf{r}}_{j}^{(v)})$ 和温度 $\tau$。
逐视图对齐损失为：
\begin{equation}
  \mathcal{L}_{\text{align}}^{(v)} =
  \frac{1}{B}\sum_{i=1}^{B}
  -\log \frac{\exp(\mathbf{S}_{ii}/\tau)}{\sum_{j=1}^{B}\exp(\mathbf{S}_{ij}/\tau)}.
\end{equation}
我们学习视图权重 $\pi_v=\mathrm{softmax}(\boldsymbol{\phi})_v$ 并通过 $s_{\text{mv}}$ 缩放总对齐：
\begin{equation}
  \mathcal{L}_{\text{mv-align}} = s_{\text{mv}}\sum_{v=1}^{V}\pi_v \mathcal{L}_{\text{align}}^{(v)}.
\end{equation}

\paragraph{冷启动重加权}
为避免对齐被高频物品主导，我们根据物品流行度 $\mathrm{pop}(i)$ 对样本进行重加权：
\begin{equation}
  w_i = 1 + \beta\cdot \max\Big(0, \frac{T-\mathrm{pop}(i)}{T}\Big),
\end{equation}
并在相似度矩阵上计算加权交叉熵（由 $\sum_i w_i$ 归一化）。

\subsection{训练目标}
我们优化：
\begin{equation}
  \mathcal{L} = \mathcal{L}_{\text{rec}} + \lambda \cdot \mathcal{L}_{\text{mv-align}},
\end{equation}
其中 $\mathcal{L}_{\text{rec}}$ 是标准的序列推荐损失（训练中对采样的正负样本进行交叉熵），$\lambda$ 是对齐权重。

\subsection{模型复杂度}
我们分析每个模型变体的空间（参数）和时间复杂度，遵循~\cite{kang2018sasrec}中的分析。
设 $|I|$ 表示物品数量，$L$ 表示最大序列长度，$d$ 表示隐藏维度，$B$ 表示批次大小，$d_t$ 表示 \tfidf 特征维度，$d_l$ 表示 \llm 嵌入维度，$V$ 表示文本视图数量，$d_v$ 表示每视图 \llm 维度，$C$ 表示交叉层数量。

\paragraph{空间复杂度（参数）}
表~\ref{tab:complexity} 总结了每种配置的参数和时间复杂度：

\begin{table*}[t]
\caption{模型变体的参数和时间复杂度。$|I|$：物品数，$L$：序列长度，$d$：隐藏维度，$d_t$：\tfidf 维度，$d_l$：\llm 维度，$V$：视图数，$d_v$：每视图维度。}
\label{tab:complexity}
\centering
\small
\begin{tabular}{lll}
\toprule
\textbf{模型} & \textbf{空间复杂度（参数）} & \textbf{时间复杂度（每批次）} \\
\midrule
\base（仅ID） & $O(|I|d + Ld + d^2)$ & $O(L^2d + Ld^2 + |I|d)$ \\
\base + \tfidf & $O(|I|d + Ld + d^2 + d_t d)$ & $O(L^2d + Ld^2 + |I|d_t d + |I|d)$ \\
\base + \tfidf + \llm（单视图） & $O(|I|d + Ld + d^2 + (d_t{+}d_l)d)$ & $O(L^2d + Ld^2 + |I|(d_t{+}d_l)d + |I|d)$ \\
\model（多视图） & $O(|I|d + Ld + d^2 + Vd_v d + d_t d)$ & $O(L^2d + Ld^2 + V|I|d_v d + |I|d)$ \\
\bottomrule
\end{tabular}
\end{table*}

\paragraph{空间分解}
\base 骨干贡献 $O(|I|d)$ 用于物品嵌入，$O(Ld)$ 用于位置嵌入，$O(d^2)$ 用于 Transformer 编码器参数（自注意力 $Q,K,V,O$ 投影和逐点 FFN 层；层数是一个小常数）。
文本变体添加与文本特征维度成比例的投影层：\tfidf 添加 $O(d_t d)$，单视图 \llm 添加 $O((d_t + d_l)d)$，多视图添加 $O(V d_v d)$ 用于逐视图投影和 SENet 模块，加上 $O(d_t d)$ 用于基础视图。
交叉层每层贡献 $O(Cd^2)$，但由于 $C$ 是一个小常数（通常为1），我们将其吸收到 $O(d^2)$ 项中。

\paragraph{时间分解}
Transformer 自注意力层花费 $O(L^2d)$ 计算 $L\times L$ 注意力矩阵，加上 $O(Ld^2)$ 用于逐点 FFN。
全排序打分增加 $O(|I|d)$ 用于计算所有物品的 logits。
文本模型产生额外的融合开销：
\begin{itemize}[leftmargin=*,itemsep=1pt]
  \item \textbf{\tfidf}：文本投影 $O(|I| \cdot d_t \cdot d)$ 和所有物品的交叉网络融合 $O(|I| \cdot d^2)$。
  \item \textbf{单视图}：由于输入维度更大，投影成本增加到 $O(|I| \cdot (d_t + d_l) \cdot d)$。
  \item \textbf{多视图}：逐视图 SENet 和投影 $O(V \cdot |I| \cdot d_v \cdot d)$ 加上拼接投影 $O(|I| \cdot Vd \cdot d)$。
\end{itemize}
对齐损失为每个视图增加 $O(B^2 d)$ 用于 InfoNCE 相似度矩阵计算。
注意 \llm 嵌入是\emph{离线预计算}并作为静态缓冲区存储的；在线推理成本主要由投影和交叉层决定，而非 \llm 编码器。

\section{两阶段训练协议}
我们使用在训练脚本中实现的两阶段训练协议：
\begin{itemize}[leftmargin=*,itemsep=2pt]
  \item \textbf{预热（可选）} 在禁用文本/对齐的情况下训练仅ID骨干几个 epoch。
  \item \textbf{A阶段（文本预热）} 冻结骨干并训练文本头和交叉模块；通过验证指标网格搜索 $(\lambda,\tau)$。
  \item \textbf{B阶段（联合微调）} 解冻骨干并使用分组学习率（文本头、交叉/DNN、骨干）微调所有参数。
\end{itemize}

\begin{algorithm}[t]
\caption{两阶段训练（预热 + A阶段 + B阶段）}
\label{alg:two-phase}
\begin{algorithmic}[1]
\Require 数据集 $\mathcal{D}$；骨干 $\base$；文本特征；网格 $\Lambda,\mathcal{T}$
\State \textbf{（可选）预热：}禁用文本训练 $\base$
\State \textbf{A阶段：}冻结骨干参数；对 $(\lambda,\tau)\in \Lambda\times\mathcal{T}$ 训练并在验证集上选择最佳
\State \textbf{B阶段：}加载最佳A阶段检查点；解冻骨干；使用分组学习率微调
\end{algorithmic}
\end{algorithm}

\section{实验}
\subsection{数据集和评估}
我们在两个亚马逊领域进行评估：\textbf{美妆}和\textbf{玩具游戏}。
我们遵循时间顺序划分，并报告 $K\in\{5,10,20\}$ 下的\textbf{全排序}指标。
除了标准指标（召回率、NDCG、MRR、命中率、精确率），我们还报告按物品交互次数的\textbf{流行度分层}指标：
\begin{itemize}[leftmargin=*,itemsep=1pt]
  \item \textbf{new}：$[1,3)$ 次交互
  \item \textbf{few}：$[3,10)$ 次交互
  \item \textbf{frequent}：$[10,\infty)$ 次交互
\end{itemize}

\subsection{基线}
我们比较四种主要配置（对应我们的运行脚本）：
\begin{itemize}[leftmargin=*,itemsep=2pt]
  \item \textbf{\base（仅ID）} 不带文本/交叉/对齐的纯序列基线。（\texttt{run50epBase\_stratified.sh}）
  \item \textbf{\base + \tfidf} 使用预计算 \tfidf 嵌入、显式交叉和对齐的轻量级文本基线。（\texttt{two\_phase\_run\_tfidf\_stratified.sh}）
  \item \textbf{\base + \tfidf + \llm（单视图）} 添加与 \tfidf 拼接的单个 \llm 嵌入（Qwen2.5-7B），使用 SENet 门控和对齐。（\texttt{two\_phase\_run\_tfidf\_llm\_stratified.sh}）
  \item \textbf{\model（多视图）} \tfidf 基础视图 + 4个多视图 \llm 提示嵌入（身份/功能/受众/类别），带逐视图 SENet、显式交叉和多视图对比对齐。（\texttt{two\_phase\_run\_multiview\_v2\_stratified.sh}）
\end{itemize}
我们还包括代表性的序列基线（如 BERT4Rec~\cite{sun2019bert4rec}、S3-Rec~\cite{zhou2020s3rec}）和轻量级文本编码器作为额外比较。

\subsection{实现细节}
除非另有说明，我们使用嵌入维度 $d=256$，最大序列长度 $L=50$，2层 Transformer 带2个头，批次大小512，Adam 优化器基础学习率 $10^{-4}$。

\paragraph{两阶段计划}
我们遵循预热（10个 epoch）$\rightarrow$ A阶段（15个 epoch）$\rightarrow$ B阶段（40个 epoch），A阶段在美妆数据上通过 \texttt{MRR@10} 选择 $(\lambda,\tau)$，在玩具游戏上通过 \texttt{Recall@10}。
我们使用分组学习率：$\text{lr}_{\text{text}}=1\times10^{-3}$ 用于文本头，$\text{lr}_{\text{cross}}=5\times10^{-4}$ 用于交叉/DNN 块，$\text{lr}_{\text{backbone}}=\text{lr}_{\text{text}}\times 0.1$ 在B阶段。

\paragraph{网格}
对于主要实验，我们在美妆上使用 $\lambda\in\{0.03\}$ 和 $\tau\in\{0.1\}$，在玩具游戏上使用 $\lambda\in\{0.05\}$ 和 $\tau\in\{0.05\}$；可通过训练脚本启用额外网格。

\paragraph{可复现性}
我们使用种子2025，报告 $K\in\{5,10,20\}$ 的全排序指标，流行度分层阈值为 \texttt{new}:[1,3)，\texttt{few}:[3,10)，\texttt{frequent}:[10,$\infty$)。
所有精确设置可从仓库提供的YAML配置和运行脚本复现。

\paragraph{参数数量}
表~\ref{tab:params} 显示了亚马逊美妆数据集上的实际可训练参数数量（$|I|=12,102$，$d=256$，$d_t=256$，$d_l=1024$，$V=4$，$d_v=1024$）。
多视图模型相比单视图基线仅增加适度的开销（约6M参数），而 \llm 嵌入本身是离线预计算并作为静态缓冲区存储的（不计入可训练参数）。

\begin{table}[h]
\caption{可训练参数数量（亚马逊美妆）}
\label{tab:params}
\centering
\small
\begin{tabular}{lrr}
\toprule
\textbf{模型} & \textbf{参数量} & \textbf{相对基线开销} \\
\midrule
\base（仅ID） & ${\sim}3.5$M & --- \\
\base + \tfidf & ${\sim}4.5$M & +1.0M \\
\base + \tfidf + \llm（单视图） & ${\sim}5.3$M & +1.8M \\
\model（多视图，$V{=}4$） & ${\sim}11.5$M & +8.0M \\
\bottomrule
\end{tabular}
\end{table}

\section{主要结果}
表~\ref{tab:main} 报告了美妆和玩具游戏的全排序结果。
在两个数据集上，添加文本特征相比仅ID基线持续改善，其中 \tfidf 提供了显著增益（美妆上 Recall@10 提升 +27.5\%，玩具上 +14.9\%）。
主要发现如下：

\paragraph{文本特征相比仅ID带来大幅增益}
在美妆上，仅 \tfidf 就将 Recall@10 从 4.69\% 提升到 5.98\%（+27.5\%），NDCG@10 从 2.69\% 提升到 3.79\%（+40.9\%）。
在玩具上，改进分别为 +14.9\% Recall@10 和 +33.1\% NDCG@10。

\paragraph{多视图呈现数据集依赖的缩放}
在\textbf{美妆}上，\model（14B）取得最佳结果，Recall@10=6.22\%（相比仅ID +32.6\%），但32B\emph{并未}比14B更好，表明在固定嵌入预算下收益递减。
在\textbf{玩具}上，多视图收益更一致地缩放：32B 达到 Recall@10=7.02\%（相比仅ID +17.6\%），在所有分层上优于较小模型。

\paragraph{冷启动（new/few）改进}
多视图在\textbf{new}和\textbf{few}分层上提供最大的相对增益。
在美妆（14B）上，Recall\_new@10 从 1.71\% 提升到 2.27\%（+32.7\%），Recall\_few@10 从 3.33\% 提升到 4.05\%（+21.6\%）。
在玩具（32B）上，Recall\_new@10 提升 +24.1\%，Recall\_few@10 提升 +20.4\%。

\begin{table*}[t]
\caption{全排序下的主要结果。每个数据集的最佳结果用\textbf{粗体}标注。所有数值为百分比（\%）。}
\label{tab:main}
\centering
\small
\begin{tabular}{lcccccccc}
\toprule
\multirow{2}{*}{\textbf{模型}} &
\multicolumn{2}{c}{\textbf{整体}} &
\multicolumn{2}{c}{\textbf{new}} &
\multicolumn{2}{c}{\textbf{few}} &
\multicolumn{2}{c}{\textbf{frequent}} \\
\cmidrule(lr){2-3}\cmidrule(lr){4-5}\cmidrule(lr){6-7}\cmidrule(lr){8-9}
& Recall@10 & NDCG@10 & Recall@10 & NDCG@10 & Recall@10 & NDCG@10 & Recall@10 & NDCG@10 \\
\midrule
\multicolumn{9}{l}{\textbf{亚马逊美妆}} \\
\base（仅ID） & 4.69 & 2.69 & 1.71 & 1.01 & 3.33 & 1.98 & 7.86 & 4.47 \\
\base + \tfidf & 5.98 & 3.79 & 2.08 & 1.36 & 3.78 & 2.53 & 10.27 & 6.44 \\
\base + \tfidf + \llm（7B） & 6.00 & 3.80 & 2.03 & 1.32 & 3.84 & 2.57 & 10.29 & 6.45 \\
\model（7B） & 6.05 & 3.82 & 2.09 & 1.33 & 3.85 & 2.56 & 10.37 & 6.49 \\
\model（14B） & \textbf{6.22} & \textbf{3.86} & \textbf{2.27} & \textbf{1.39} & \textbf{4.05} & \textbf{2.66} & \textbf{10.60} & \textbf{6.51} \\
\model（32B） & 6.05 & 3.81 & 2.08 & 1.34 & 3.83 & 2.53 & 10.40 & 6.47 \\
\midrule
\multicolumn{9}{l}{\textbf{亚马逊玩具游戏}} \\
\base（仅ID） & 5.97 & 3.32 & 1.91 & 1.10 & 4.61 & 2.53 & 10.61 & 5.92 \\
\base + \tfidf & 6.86 & 4.42 & 2.24 & 1.52 & 5.46 & 3.54 & 12.11 & 7.75 \\
\base + \tfidf + \llm（7B） & 6.83 & 4.39 & 2.26 & 1.51 & 5.31 & 3.46 & 12.10 & 7.73 \\
\model（7B） & 6.90 & 4.43 & 2.30 & 1.53 & 5.40 & 3.51 & 12.21 & 7.79 \\
\model（14B） & 6.95 & 4.43 & 2.33 & 1.52 & 5.47 & 3.56 & 12.26 & 7.78 \\
\model（32B） & \textbf{7.02} & \textbf{4.47} & \textbf{2.37} & \textbf{1.56} & \textbf{5.55} & \textbf{3.60} & \textbf{12.38} & \textbf{7.81} \\
\bottomrule
\end{tabular}
\end{table*}

表~\ref{tab:mrr} 报告了用于排序质量分析的 MRR 指标。

\begin{table}[t]
\caption{MRR@10 结果（\%）。$\Delta$ = 相对仅ID的相对提升。}
\label{tab:mrr}
\centering
\small
\begin{tabular}{lcccc}
\toprule
\textbf{模型} & \textbf{MRR} & $\Delta$ & \textbf{MRR\_new} & \textbf{MRR\_few} \\
\midrule
\multicolumn{5}{l}{\textbf{亚马逊美妆}} \\
\base（仅ID） & 2.07 & --- & 0.79 & 1.56 \\
\base + \tfidf & 3.12 & +50.7\% & 1.13 & 2.14 \\
+ \llm（7B） & 3.13 & +51.2\% & 1.10 & 2.18 \\
\model（14B） & \textbf{3.13} & \textbf{+51.2\%} & \textbf{1.12} & \textbf{2.23} \\
\midrule
\multicolumn{5}{l}{\textbf{亚马逊玩具游戏}} \\
\base（仅ID） & 2.49 & --- & 0.84 & 1.87 \\
\base + \tfidf & 3.66 & +47.0\% & 1.28 & 2.95 \\
+ \llm（7B） & 3.64 & +46.2\% & 1.27 & 2.90 \\
\model（32B） & \textbf{3.68} & \textbf{+47.8\%} & \textbf{1.31} & \textbf{3.00} \\
\bottomrule
\end{tabular}
\end{table}

\section{分析}
\subsection{消融实验}
我们消融 (i) 显式交叉，(ii) 对齐，(iii) 逐视图归一化，和 (iv) 冷启动重加权。
我们还研究文本视图数量和多视图对齐缩放 $s_{\text{mv}}$ 的影响。

\paragraph{多视图增益的驱动因素}
由于多视图文本是 \model 的核心，我们消融其两个核心成分：显式交叉融合和逐视图 SENet 门控。
移除交叉模块可以略微增加 Recall@K，但会急剧降低排序质量（如在美妆上 MRR@10 相对下降约20\%，NDCG@10 下降约10\%），表明交叉交互对于在全排序下将多视图召回转化为头部排名增益至关重要。
相比之下，SENet 对整体 Recall/MRR/NDCG 的影响较小，但实质性地改变了曝光：禁用 SENet 大幅增加 Coverage\_new，这与校准门控偏向头部物品而牺牲冷启动覆盖一致。
综合来看，这些结果表明多视图特征的有效性主要取决于视图\emph{如何}融合（交叉）和校准（SENet），而不仅仅是添加更多文本信号。

\paragraph{何时使用多视图 vs 单视图}
我们的结果提供了一个简单的经验法则。
当目标强调 Recall@K（特别是 $K{=}10/20$）和\textbf{few}/\textbf{frequent}分层性能时，多视图提示最有价值，此时互补的提示视图提供更广的语义覆盖，可以改善候选召回。
相反，当大部分增益已经来自 \tfidf{+}\llm（单视图）时——如在固定嵌入预算下的美妆——或当主要目标是冷启动排序和曝光（Coverage\_new 和\textbf{new}分层 MRR/NDCG）时，单视图 \llm 基线可能是更稳健和性价比更高的选择。

\subsection{实证要点：文本加权和预算缩放}
我们总结三个在美妆和玩具游戏上一致的观察，有助于理解文本增强序列推荐器的权衡。

\textbf{(A) 文本加权以较小的召回变化改善排序质量：}
增加文本分支的有效贡献（如更强的文本融合权重和更温和的门控正则化）往往只在 Recall@K 上产生小幅波动，同时持续改善排序导向的指标如 MRR 和 NDCG。

\textbf{(B) LLM 缩放在部署约束下不是单调可靠的：}
更大的 \llm 编码器不一定转化为更好的排序或冷启动指标，表明朴素的"缩放法则"在此设置下不成立。

\textbf{(C) 多视图增益是数据依赖的，可能降低冷启动/新物品覆盖：}
虽然多视图提示可以改善整体或高频物品性能，但净收益因领域而异，我们观察到\textbf{new}物品覆盖率降低的趋势，表明存在头部-尾部权衡。

\paragraph{跨领域综合}
在两个领域中，最稳健的改进来自用单个 \llm 视图丰富 \tfidf，而多视图提示并不一致地优于单视图基线，可能将增益转向头部（高频）物品，代价是冷启动覆盖（Coverage\_new）和\textbf{new}分层排序指标。
此外，扩大 \llm 规模最多只能获得分层和指标特定的增益（通常在高频分层上最明显），并不能转化为普遍更好的 MRR/NDCG 或冷启动性能。

\paragraph{关于嵌入预算和压缩的说明}
重要的是，我们的多视图 \llm 特征是在固定嵌入预算下生成的以保证实用性：每个视图在融合前通过 TruncatedSVD 压缩到低维向量（如64维）。
这种压缩作为表示瓶颈，是分析 \llm 缩放时的关键混淆因素；更大的编码器可能产生更高秩的表示，在相同预算下被更激进地截断，可能掩盖其优势。
因此，我们将缩放行为解释为给定压缩强度下的\emph{预算缩放}，并将压缩预算的系统研究（如变化 SVD 维度和跟踪解释方差保留）作为未来的消融实验。

\subsubsection{固定嵌入预算下的领域观察}
在固定的多视图嵌入预算下（每视图 TruncatedSVD 到64维），美妆主要由 \tfidf$\rightarrow$\tfidf{+}\llm 驱动，而多视图变体相对于单视图 \llm 基线几乎没有额外的整体改进，并且在冷启动（\textbf{new}）分层上持续表现不佳——包括 Recall\_new/NDCG\_new/MRR\_new 和 Coverage\_new——在32B时降级最大（表~\ref{tab:main}）。
相反，在玩具游戏上，多视图提示相对于单视图 \llm 基线改善了整体 Recall@10/20 和 few/frequent 分层，并在 frequent 分层上从 7B$\rightarrow$14B$\rightarrow$32B 显示出弱缩放趋势；然而，MRR/NDCG 不单调缩放，Coverage\_new 在更大规模时可能下降，揭示了在32B时最明显的召回-排序/覆盖权衡。

\subsection{采样指标 vs 全排序指标}
我们比较采样评估（如 uni100）和全排序，并量化不同设置下的模型排序不一致性。
我们测试对齐是否减少这种不匹配并提高鲁棒性。

\subsection{效率和权衡}
我们报告 GPU 内存占用、吞吐量和推理延迟。
由于 \llm 嵌入是离线预计算的，在线成本主要由投影/交叉层和特征存储决定，在节约型（仅 \tfidf）和更强（多视图）选项之间实现清晰的权衡。

\subsection{可解释性}
我们可视化 (i) 视图门控 $g_v$，(ii) 学习的视图对齐权重 $\pi_v$，和 (iii) 跨视图和ID空间的表示几何（t-SNE/UMAP）。
我们还检查交叉模块贡献（如 top-$k$ 交叉权重）以理解哪些交互驱动长尾增益。

\begin{figure}[t]
\centering
\IfFileExists{figures/framework.pdf}{%
  \includegraphics[width=\linewidth]{figures/framework.pdf}%
}{%
  \fbox{\parbox[c][0.22\textheight][c]{0.95\linewidth}{\centering
  \textbf{框架图占位符}\\
  请将 framework.pdf 放在 paper\_sigir/figures/ 下。}}%
}
\caption{框架概览（占位符）}
\label{fig:framework}
\end{figure}

\section{相关工作}
\paragraph{序列推荐}
基于 Transformer 的序列模型如 \base~\cite{kang2018sasrec} 和 BERT 风格目标~\cite{sun2019bert4rec} 是强大的骨干。

\paragraph{文本增强推荐}
物品文本特征（来自词袋、神经编码器或 \llm 嵌入）提供补充信号；我们的重点是多视图提示和可部署的融合。

\paragraph{对比对齐和评估}
我们基于对比目标（\InfoNCE~\cite{oord2018cpc}）构建，并研究采样评估如何误导模型选择~\cite{krichene2020sampledmetrics}。

\section{结论}
我们提出了 \model，一个结合归一化多视图文本特征、显式交叉交互和多视图对比对齐的序列推荐插件框架。
我们在亚马逊美妆和玩具游戏上的实验揭示了几个实用洞察：

\textbf{(1) 文本特征带来显著增益}：仅 \tfidf 就在美妆上提升 Recall@10 +27.5\%，在玩具上提升 +14.9\%，相比仅ID基线，NDCG 增益更大（+40.9\% 和 +33.1\%）。

\textbf{(2) 多视图缩放因数据集而异}：在美妆上，14B 取得最佳结果（相比仅ID Recall@10 +32.6\%），但 32B 未能进一步提升。在玩具上，32B 持续优于较小模型（Recall@10 +17.6\%）。

\textbf{(3) 冷启动收益一致}：多视图在\textbf{new}和\textbf{few}分层上提供最大的相对增益，美妆（14B）上 Recall\_new@10 提升 +32.7\%，玩具（32B）上提升 +24.1\%。

\textbf{(4) 预算缩放有限制}：在固定嵌入压缩下，更大的 \llm 编码器并不保证单调改进，突显了领域特定评估的必要性。

\balance  % 最后一页平衡列
\bibliographystyle{ACM-Reference-Format}
\bibliography{refs}

\end{document}
